%==============================================================================
\section{Statistics --- Thống kê trong học máy}
\label{sec:statistics}
%==============================================================================

\begin{dongco}[title={Vai trò của thống kê}]
Thống kê là công cụ then chốt trong học máy vì giúp hiểu pattern ẩn trong dữ liệu.
Nó cung cấp phương pháp mô tả, tóm tắt và phân tích dữ liệu.
\end{dongco}

\subsection{Statistics là gì?}

\begin{dinhnghia}[title={Định nghĩa}]
Thống kê là nhánh toán học về thu thập, phân tích, diễn giải và trình bày dữ liệu.
Nó cung cấp các phương pháp và kỹ thuật phân tích dữ liệu và rút kết luận.
\end{dinhnghia}

\begin{ynghia}[title={Nền tảng cho ML}]
Thống kê là nền tảng học máy: phân tích và trực quan hoá để tìm pattern.
Ứng dụng: xác thực mô hình, làm sạch dữ liệu, chọn mô hình, đánh giá hiệu năng, \ldots
\end{ynghia}

\subsection{Các khái niệm thống kê cơ bản cho ML}

\begin{phuongphap}[title={Danh sách khái niệm quan trọng}]
\begin{itemize}
  \item \textbf{Trung bình (Mean), Trung vị (Median), Mode}: đo xu hướng trung tâm.
  \item \textbf{Phương sai (Variance), Độ lệch chuẩn (Standard deviation)}: độ phân tán quanh mean.
  \item \textbf{Phân vị (Percentiles)}: giá trị dưới đó một tỷ lệ \% quan sát rơi vào.
  \item \textbf{Phân phối dữ liệu (Data distribution)}: cách điểm dữ liệu trải trên miền giá trị.
  \item \textbf{Độ lệch (Skewness), Độ nhọn (Kurtosis)}: độ lệch và độ nhọn của phân phối.
  \item \textbf{Sai số (Bias), Phương sai (Variance)}: nguồn sai số dự đoán mô hình.
  \item \textbf{Giả thuyết (Hypothesis)}: giả thuyết / lời giải đề xuất cho bài toán.
  \item \textbf{Hồi quy tuyến tính (Linear Regression)}: dự đoán giá trị của một biến dựa trên giá trị của một biến khác .
  \item \textbf{Hồi quy logistic (Logistic Regression)}: ước tính xác suất xảy ra của một sự kiện.
  \item \textbf{Phân tích thành phần chính (PCA)}: giảm số chiều của dữ liệu bằng cách tìm ra các thành phần chính (principal components) của dữ liệu.
\end{itemize}
\end{phuongphap}

\subsection{Hai loại thống kê}

\begin{dinhnghia}[title={Thống kê mô tả}]
Tập quy tắc/phương pháp dùng để \textbf{mô tả hoặc tóm tắt} đặc trưng của dataset.
\end{dinhnghia}

\begin{dinhnghia}[title={Thống kê suy luận}]
Đối phó với \textbf{dự đoán và suy luận} về quần thể dựa trên mẫu dữ liệu.
\end{dinhnghia}

\subsubsection{Thống kê mô tả (Descriptive Statistics)}

\begin{ynghia}[title={Nội dung}]
Nhánh thống kê tóm tắt và phân tích dữ liệu: mean, median, mode, variance, standard deviation.
Giúp hiểu xu hướng trung tâm, độ biến thiên và hình dạng phân phối.
\end{ynghia}

\begin{phuongphap}[title={Ứng dụng trong ML}]
Tóm tắt dữ liệu, nhận diện outlier, phát hiện pattern.
Ví dụ: dùng mean và standard deviation để mô tả phân phối một feature.
\end{phuongphap}

\begin{vidu}[title={Tính thống kê mô tả với NumPy và Pandas}]
\begin{lstlisting}
import numpy as np
import pandas as pd

data = np.array([1, 2, 3, 4, 5])
df = pd.DataFrame(data, columns=["Values"])
print(df.describe())
\end{lstlisting}

\textbf{Output} (tóm tắt gồm count, mean, std, min, các phân vị, max):
\begin{lstlisting}
         Values
count    5.000000
mean     3.000000
std      1.581139
min      1.000000
25%      2.000000
50%      3.000000
75%      4.000000
max      5.000000
\end{lstlisting}
\end{vidu}

\subsubsection{Thống kê suy luận (Inferential Statistics)}

\begin{ynghia}[title={Nội dung}]
Nhánh thống kê dự đoán và suy luận về quần thể từ mẫu: kiểm định giả thuyết, khoảng tin cậy, hồi quy.
\end{ynghia}

\begin{phuongphap}[title={Ứng dụng trong ML}]
Dự đoán trên dữ liệu mới dựa trên dữ liệu hiện có --- ví dụ hồi quy giá nhà theo số phòng, diện tích, \ldots
\end{phuongphap}

\begin{vidu}[title={Hồi quy OLS với statsmodels}]
\begin{lstlisting}
import statsmodels.api as sm
import numpy as np

X = np.array([1, 2, 3, 4, 5])
y = np.array([2, 4, 6, 8, 10])

X = sm.add_constant(X)
model = sm.OLS(y, X).fit()

print(model.summary())
\end{lstlisting}

\textbf{Output} (trích): bảng hệ số, std err, $t$, $P>|t|$, $R^2$; với dữ liệu tuyến tính hoàn hảo, hệ số $x_1 \approx 2$.
\end{vidu}
